\section{Building nets for spaces of functions}
\label{sec:lipsort}

The repository's \repo{docs/proofs/foundations/ipra_finite_use_framework_2026-11-15/sections/09_compactness.tex}{finite-net transfer theorem} starts with a supplied net. We construct such a net for a function space and compute the pseudometric on its codes. This produces the input for a second application of the same construction to functionals.

Fix a nonempty rational pseudometric presentation $(D,d)$ as in Section~\ref{sec:base}, with its dyadic nets. Write $X$ for the completion and fix a nonnegative rational $L$. The target sort is the space $\Lip_L(X,[0,1])$ of functions satisfying $|f(x)-f(y)|\le Ld(x,y)$, with the uniform distance. A finite code $A$ is a list of pairs $(a,b)$ with $a\in D$ and $b\in\Q\cap[0,1]$. Its value at a rational-carrier point is
\begin{equation}
 G_A(x)=\min\bigl(\{1\}\cup\{b+Ld(x,a):(a,b)\in A\}\bigr).
 \label{eq:cone}
\end{equation}
The empty list represents the constant one. This uses the upper extension formula of McShane for Lipschitz data. We allow mutually inconsistent sample heights and do not require $G_A(a)=b$. This detail makes the enumeration and exact distance calculation work together.\footnote{E. J. McShane, \emph{Extension of range of functions}, Bull. Amer. Math. Soc. 40 (1934), 837--842. \href{https://doi.org/10.1090/S0002-9904-1934-05978-0}{Original article}.}

\begin{lemma}[Distance witnessed by the anchors]\label{lem:anchors}
Every $G_A$ is $[0,1]$-valued and $L$-Lipschitz. Set
\[
 E(A,h)=\max\bigl(\{0\}\cup\{h(a)-b:(a,b)\in A\}\bigr).
\]
For any $[0,1]$-valued $L$-Lipschitz $h$, and $e\ge0$,
\[
  E(A,h)\le e\quad\Longleftrightarrow\quad
      h(x)-G_A(x)\le e\quad\text{for every }x.
\]
In particular, the number
\begin{equation}
 \Delta(A,B)=\max\{E(A,G_B),E(B,G_A)\}                 \label{eq:metric}
\end{equation}
is the exact uniform distance of two finite codes.
\end{lemma}
\begin{proof}
Triangle inequalities show each cone $b+Ld(x,a)$ changes by at most $Ld(x,y)$. A finite minimum of such functions obeys the same bound: every entry at $x$ is at least the corresponding entry at $y$ minus $Ld(x,y)$, and conversely. The constant entry $1$ caps the range, and the heights and distances are nonnegative.

If $E(A,h)\le e$, then for each anchor $h(x)\le h(a)+Ld(x,a)\le b+Ld(x,a)+e$. Also $h(x)\le 1+e$. Minimize the right sides. In the converse direction use $x=a$ and $G_A(a)\le b$. Applying both directions with $h=G_B$ and then $h=G_A$ proves that $\Delta$ bounds all discrepancies and is below every nonnegative such bound. If $\Delta>0$, a maximizing term in one of the lists supplies an actual point of attainment. If $\Delta=0$, all values agree. Symmetry and reflexivity follow directly; apply the pointwise triangle inequality and the bound characterization to get the triangle law for $\Delta$. All operations in \eqref{eq:metric} are rational, including when two different codes have distance zero.
\end{proof}

\begin{theorem}[A primitive recursive inner net]\label{thm:inner}
Let $D_r$ have $N(r)$ entries and radius $\eta=2^{-r}$. For a requested $k$, put
\[
 h=2^{-k-2},\qquad
 r=k+3+\min\{c:2^c\ge\max(1,L)\}.
\]
On the list $D_r$ take all height vectors with entries in $\{0,h,2h,\ldots,1\}$. The resulting $(2^{k+2}+1)^{N(r)}$ codes are members of the Lipschitz ball and cover it in uniform distance with radius at most $2^{-k-1}$. A bounded procedure selects a covering code from the values of its input function at the listed points, each read to absolute error $h/2$.
\end{theorem}
\begin{proof}
Given $f$, round its rational sample at $a$ to the nearest multiple of $h$, resolving exact rational ties towards the lower grid index; clip first to $[0,1]$. The resulting $b$ differs from $f(a)$ by at most $\delta=h$. The lower bound $G_A(x)\ge f(x)-\delta$ follows from Lemma~\ref{lem:anchors}. For the upper bound, choose an anchor $a$ within $\eta$ of $x$ on the dense carrier. Then
\[
 G_A(x)\le b+Ld(x,a)
     \le f(a)+\delta+L\eta
     \le f(x)+\delta+2L\eta.
\]
Our choice of $r$ gives $2L\eta\le h$, hence the two-sided error is at most $2h=2^{-k-1}$. On named inputs the same bound follows by approximating $x$ with dense-carrier elements to arbitrary precision; both functions have modulus $L$. For named $f(a)$, each occurrence of a value comparison can be written with independent rational errors as in Section~\ref{sec:base}. The height choice uses one finite read. In subsequent proofs of its bound one may take finer reads of the \emph{same} $f$ and compare by regularity. No comparison of a real value with a grid separator is needed.

The exponent $c$ can be found by scanning up to $1+\lceil\max(1,L)\rceil$. Encode the selected vector as a radix-$B$ numeral, $B=2^{k+2}+1$, with a fixed digit order and $N(r)$ positions, including leading zeros. Its code is strictly below the displayed count. The evaluation loop takes a prefix minimum including $1$; comparison to a proposed lower bound is the bounded conjunction of the corresponding prefix inequalities. The distance loop takes a maximum including $0$ over each anchor list. Testing selected heights and membership of anchors uses bounded bits. These observations give QF induction invariants for each loop even when $N(r)$ varies. The input net supplies the only spatial coverage premise.
\end{proof}

\paragraph{Why the completion is the intended function space.}
For a regular sequence of valid codes $A_j$, $G_{A_j}(x)$ is a regular scalar sequence at each dense $x$ by \eqref{eq:metric}. Its limit $f(x)$ lies in $[0,1]$ and has the same Lipschitz bound: apply the bound for $G_{A_j}$ and remove two vanishing evaluation margins. Moreover $\lvert f(x)-G_{A_j}(x)\rvert\le2^{-j}$ uniformly in $x$ by comparison with later codes. Extend this $f$ to $X$ by the explicit sampling schedule in Section~\ref{sec:base}. For a name $x_i$, for example, read $A_{k+2}$ and $x_{k+2+c}$, where $2^c\ge\max(1,L)$, and evaluate the rational expression \eqref{eq:cone}. The two errors sum to at most $2^{-k-1}$. Two equivalent sequences yield the same function by the same bound. Conversely, apply Theorem~\ref{thm:inner} at every $k$ to a given $f$. The uniform approximation bounds yield $\Delta(A_k,A_l)\le2^{-k-1}+2^{-l-1}$ and recover $f$ on evaluation. For presented functions the selector is PR jointly in all numerical parameters whenever their sample procedure is. Externally the same density proof in ordinary mathematics identifies the metric completion with the full classical Lipschitz ball. In this external reading one may round classical values or use the existence of finite rational approximations. The constructive version uses the specified rational observations.

\paragraph{Higher-type expressions with bounded scalar range.}
Starting from nonempty net-equipped rational pseudometric sorts, close under finite products and $X\mapsto\Lip_L(X,[0,1])$ for rational $L\ge0$. At each function constructor the dense carrier consists of the \emph{valid} codes, with \eqref{eq:metric}; here validity includes the height bounds. Equations \eqref{eq:cone} and \eqref{eq:metric} use only operations already supplied at the preceding sort. The constructor proves their preconditions for the next sort. Evaluation is jointly continuous with
\[
 |f(x)-g(y)|\le d_\infty(f,g)+Ld(x,y).
\]
Consequently $1+L$ is a Lipschitz bound for evaluation on the max product. Structural recursion in the metatheory supplies a PR evaluator, modulus, net and coverage proof for every fixed sort expression. Combining this with Lemma~\ref{lem:sup} gives arithmetic values for nested $\inf,\sup$ over functions and functionals for observables with joint evaluators and certified moduli. For example, a real parameter $f$ is an input point of $Y=\Lip_1([0,1],[0,1])$; the map $f\mapsto f(1/2)$ is a point of $\Lip_1(Y,[0,1])$. At this second level, the inner distances in \eqref{eq:cone} are computed by \eqref{eq:metric}, so no recursive call to an oracle for the supremum of two functions is necessary.

There is also a modulus version. If $\omega$ is a nonnegative, nondecreasing, subadditive rational PR function on nonnegative rationals, with $\omega(0)=0$ and a PR rate at zero, replace $d$ by $\omega\circ d$. The triangle inequality follows first by monotonicity and then by subadditivity. The rate selects finer $d$-nets giving dyadic nets in the new distance. The $1$-Lipschitz construction then covers precisely functions satisfying $|f(x)-f(y)|\le\omega(d(x,y))$ on the original carrier. A non-injective change of pseudometric merely identifies codes at distance zero.

\paragraph{Two boundaries with proofs.}
The space of \emph{all} continuous $[0,1]$-valued functions on $[0,1]$ has no uniform finite $1/3$-net. On each $[2^{-j-1},2^{-j}]$ place a triangular tent of height one, zero at both ends, and zero outside. The tents have pairwise distance one at a peak. By the triangle inequality one $1/3$-ball contains at most one; one more tent than the number of balls gives the contradiction by finite counting. A prescribed common modulus is genuine information.

Even on the constants $0,1$ a uniformly computable extraction of a convergent subsequence \emph{with convergence rate} from arbitrary PR sequences would decide halting. At step $n$ take the constant $1$ exactly when a specified machine has halted within $n$ steps and $0$ otherwise. This is a single PR family. A strictly increasing subsequence with a Cauchy rate at tolerance $1/3$ is constant after its reported index. If its value there is zero the machine never halts, since any later halt would force all sufficiently late terms to one. If it is one the machine has halted. The usual diagonal program that halts exactly when its own halting decider says it loops rules out the decider. This is a metatheoretic obstruction to a uniform producer. Total boundedness and the computation of supremum \emph{values} avoid such extraction.
